\documentclass[11pt,letterpaper]{article}
\usepackage[margin=0.75in]{geometry}
\usepackage[utf8]{inputenc}
\usepackage[T1]{fontenc}
\usepackage{lmodern}
\usepackage{helvet}
\usepackage{textcomp}
\usepackage{parskip}
\usepackage{longtable}
\usepackage{booktabs}
\usepackage{array}
\usepackage{calc}
\usepackage{amssymb}
\usepackage{enumitem}
\usepackage{titlesec}
\usepackage{xcolor}
\usepackage[colorlinks=true,urlcolor=blue,linkcolor=black]{hyperref}
\renewcommand{\familydefault}{\sfdefault}
\setlist[itemize]{leftmargin=1.3em,itemsep=1pt,topsep=2pt}
\setlist[enumerate]{leftmargin=1.5em,itemsep=1pt,topsep=2pt}
\titleformat{\section}{\large\bfseries}{}{0em}{}[\vspace{0.1em}\hrule\vspace{0.2em}]
\newcolumntype{L}[1]{>{\raggedright\arraybackslash}p{#1}}
\providecommand{\tightlist}{\setlength{\itemsep}{0pt}\setlength{\parskip}{0pt}}
\setlength{\parindent}{0pt}\setlength{\parskip}{4pt}
\begin{document}
\section*{Teaching Statement --- \'{A}ngel Luis Robles-Fern\'andez}
\textbf{Oregon State University --- Assistant Professor, Integrative Biology}

My goal is to help students use artificial intelligence to move faster from a biological question to a transparent, testable analysis. Accelerated code generation is valuable when it creates time for design, validation, and biological interpretation. I want students to become responsible designers of scientific workflows, not simply users of generated code. Computation should support observation, experiments, and theory, rather than substitute for them.

\textbf{Experience and instructional approach.}
I have taught Introduction to R, served as a teaching assistant for Quantitative Ecology and Statistical Models for Biology at Arizona State University, and been invited to teach graduate-level ecological niche modeling at INECOL. As an RStudio Certified Trainer in the Tidyverse, I use live coding, worked examples, and collaborative debugging. My experience in industrial machine learning and research software development informs the practices I will bring to AI-assisted instruction.

\textbf{From ``vibe-coding'' to software design.}
I will redirect rapid, prompt-driven coding toward a design-first workflow. Before requesting code, students will specify their biological question, inputs, outputs, assumptions, and validation criteria. They will decompose analyses into small functions, then use AI to implement or refactor bounded tasks. Readable names, code review, version control, and tests against known or simulated cases will make each step inspectable. Data provenance, units, dependencies, and reproducible runs will be part of the assignment. For example, a project on performance--temperature curves can require checks on units and numerical behavior before students interpret ecological consequences. Instructor-curated reference cases will help distinguish code that runs from an analysis that answers the biological question.

\textbf{An AI-assisted teaching platform.}
I propose piloting my platform, \href{https://learning.ecoseek.org}{learning.ecoseek.org}, alongside classroom activities and notebooks. It provides AI-generated programming exercises, including debugging, implementation, and translation tasks with executable checks, and English/Spanish content. I will adapt examples to biological coursework and review generated materials for scientific accuracy. Short practice cycles will lead into peer discussion and instructor feedback; automated checks will support practice, not replace my assessment of understanding. Classroom adoption will be subject to institutional privacy and accessibility review.

\textbf{Inclusive learning and accountable AI use.}
Ungraded diagnostics, scaffolded templates, and progressively less guided tasks will accommodate varied preparation. Active learning and rotating group roles will value biological insight as well as coding experience. Students will not be assessed on access to premium AI models; equivalent practice with worked examples and peer feedback will remain available. Students will disclose AI assistance, verify sources, and explain analytical choices. Assessment will combine reproducible notebooks, design explanations, and brief code walkthroughs or debugging tasks, emphasizing scientific reasoning over polished output. Formative assessment and student feedback will guide revisions to instruction.

\textbf{Mentoring and departmental contributions.}
I will mentor undergraduate and graduate researchers through clear expectations, regular feedback, and increasing ownership of questions and software, providing support in statistics, machine learning, principles of genetics, and population genetics. I am prepared to contribute to introductory biology, ecology, evolution, and quantitative courses in the integrative biology program. These practices will equip students for independent biological research and transferable computational work in academic and non-academic careers.
\end{document}
